\section{Kataklizmik X-Işını İkilileri (Cataclysmic X-ray Binaries)}

Kataklizmik X-ışını ikilileri, bir \emph{beyaz cüce} ile genellikle düşük kütleli ($M \lesssim 1\,M_\odot$) bir anakol yıldızından oluşan, kütle aktarımı sonucu yüksek enerjili elektromanyetik ışınım üreten yakın ikili sistemlerdir. Bu sistemler, klasik X-ışını ikililerinden farklı olarak kompakt cisim olarak nötron yıldızı veya kara delik değil, beyaz cüce içerir.

Bu nedenle kataklizmik değişenler (Cataclysmic Variables, CV) ile güçlü bir fiziksel örtüşme gösterir ve özellikle manyetik türleri X-ışını bandında baskın kaynaklar hâline gelir.

\subsection{Beyaz Cüce ve Fiziksel Özellikleri}

Beyaz cüceler, düşük ve orta kütleli yıldızların evriminin son ürünü olup elektron dejenere basıncı ile dengede tutulur. Yaklaşık kütle-yarıçap ilişkisi:
\begin{equation}
R \propto M^{-1/3}
\end{equation}

Tipik değerler:
\begin{align}
M_{\mathrm{WD}} &\sim 0.6\,M_\odot \\
R_{\mathrm{WD}} &\sim 10^4\ \mathrm{km}
\end{align}

Beyaz cüceler, nötron yıldızlarına kıyasla daha zayıf kütleçekim alanına sahip olsalar da, akresyon süreçleri X-ışını üretmek için yeterince enerjiktir.

\subsection{Oluşum Mekanizması}

Kataklizmik X-ışını ikililerinin oluşumu, ikili yıldız evrimi bağlamında şu aşamalardan geçer:
\begin{enumerate}
    \item Başlangıçta iki normal yıldızdan oluşan yakın bir ikili sistem
    \item Daha büyük kütleli yıldızın kırmızı dev evresinde ortak zarf (common envelope) evresi
    \item Zarfın atılması ve çekirdeğin beyaz cüceye dönüşmesi
    \item Yörüngenin önemli ölçüde daralması
    \item Düşük kütleli eş yıldızın Roche lobunu doldurması
\end{enumerate}

Yörünge periyodu genellikle:
\begin{equation}
P_{\mathrm{orb}} \sim 1\text{--}10\ \mathrm{hour}
\end{equation}
aralığındadır.

\subsection{Kütle Aktarımı ve Roche Lobu}

Bu sistemlerde kütle aktarımı, Roche lobu taşması yoluyla gerçekleşir. Roche lobu yarıçapı Eggleton yaklaşımı ile:
\begin{equation}
\frac{R_L}{a} = \frac{0.49q^{2/3}}{0.6q^{2/3} + \ln(1+q^{1/3})}
\end{equation}
şeklinde ifade edilir. Burada $q = M_2/M_1$ kütle oranıdır.

Kütle aktarım hızı tipik olarak:
\begin{equation}
\dot{M} \sim 10^{-11} \text{--} 10^{-8}\, M_\odot\,\mathrm{yr^{-1}}
\end{equation}

\subsection{Kataklizmik X-Işını İkililerinin Türleri}

Kataklizmik X-ışını ikilileri, beyaz cücenin manyetik alanına göre sınıflandırılır.

\subsubsection{Manyetik Olmayan Sistemler (Diskli CV'ler)}

Bu sistemlerde beyaz cücenin manyetik alanı zayıftır:
\begin{equation}
B \lesssim 10^5\ \mathrm{G}
\end{equation}

Kütle, tam gelişmiş bir akresyon diski oluşturarak beyaz cüceye düşer. X-ışını üretimi çoğunlukla disk iç bölgeleri ve sınır tabakası (boundary layer) kaynaklıdır.

\subsubsection{Ara Manyetik Sistemler (Intermediate Polars, IP)}

Bu sistemlerde manyetik alan orta şiddettedir:
\begin{equation}
10^5\ \mathrm{G} \lesssim B \lesssim 10^7\ \mathrm{G}
\end{equation}

Disk iç bölgelerde kesilir ve madde manyetik alan çizgileri boyunca kutuplara yönlendirilir. Beyaz cücenin dönme periyodu, yörünge periyodundan kısadır:
\begin{equation}
P_{\mathrm{spin}} < P_{\mathrm{orb}}
\end{equation}

\subsubsection{Polarlara (AM Herculis Sistemleri)}

Bu sistemlerde beyaz cücenin manyetik alanı çok güçlüdür:
\begin{equation}
B \sim 10^7\text{--}10^8\ \mathrm{G}
\end{equation}

Akresyon diski oluşmaz. Madde doğrudan manyetik alan boyunca kutuplara akar. Beyaz cücenin dönmesi yörünge ile senkronizedir:
\begin{equation}
P_{\mathrm{spin}} = P_{\mathrm{orb}}
\end{equation}

\subsection{X-Işını Yayınım Mekanizması}

Kataklizmik X-ışını ikililerinde X-ışını üretimi, akresyon sırasında maddenin serbest düşüşüyle ilişkilidir. Birim kütle başına açığa çıkan enerji:
\begin{equation}
\Delta E \approx \frac{GM_{\mathrm{WD}}}{R_{\mathrm{WD}}}
\end{equation}

Manyetik sistemlerde, madde beyaz cücenin yüzeyine çarpmadan önce bir \emph{şok} oluşturur. Şok sonrası sıcaklık:
\begin{equation}
kT_{\mathrm{shock}} \approx \frac{3}{8}\frac{GM_{\mathrm{WD}} m_p}{R_{\mathrm{WD}}}
\end{equation}

Bu sıcaklık:
\begin{equation}
T \sim 10^7\text{--}10^8\ \mathrm{K}
\end{equation}
mertebesinde olup sert X-ışını yayınımına neden olur.

X-ışını parlaklığı yaklaşık olarak:
\begin{equation}
L_X \approx \frac{GM_{\mathrm{WD}}\dot{M}}{R_{\mathrm{WD}}}
\end{equation}
şeklinde ifade edilir.

\subsection{Astrofiziksel Önemi}

Kataklizmik X-ışını ikilileri:
\begin{itemize}
    \item Akresyon fiziğinin düşük kütleçekim rejiminde incelenmesi
    \item Manyetik alan–plazma etkileşiminin test edilmesi
    \item İkili yıldız evrimi modellerinin sınanması
    \item Tip Ia süpernova öncüllerinin anlaşılması
\end{itemize}
açılarından yüksek enerji astrofiziğinde kritik öneme sahiptir.
